\newif\ifglobal

%%%%%%%%%%%%%%%%%%%%%%%
%%% Compile to view %%%
%%%%%%%%%%%%%%%%%%%%%%%

%%%%%%%%%%%%%%%%%%%%%%%%%%%%%%%%%%%%%%%%%%%%%%%%%%%
%%% Readme:                                     %%%
%%% https://www.overleaf.com/read/bwdjvfjksvjy/ %%%
%%%%%%%%%%%%%%%%%%%%%%%%%%%%%%%%%%%%%%%%%%%%%%%%%%%

% >>> M?: marks a module setting number or important notice

% >>> M4: Set {./relative-path} to external files for this module <<<
\def \modulefiles {./37-XTWDT}
% To add an external file, use:
% (Replace '---' with actual filename)
% '\input{\modulefiles/tables/---.tex}' for tables
% '\includegraphics{\modulefiles/figures/---}' for figures
% '\subfile{\modulefiles/---}' for register files

% >>> M1: This module is in global project? <<<
% 'YES' - uncomment; 'NO' - comment out
\globaltrue

\ifglobal
    \documentclass[main\_\_EN.tex]{subfiles}

\else
    \documentclass[openany, 10pt]{book}

    % >>> M2: In ./00-shared/preamble-trm-module, also find
    %         settings for visibility of todo notes, watermark <<<
    \usepackage{preamble-trm-module}

    % >>> M3: Temporary labels in English,
    %         you can add your own labels there <<<
    \myexternaldocument{./00-shared/temp-labels-en}

\fi %\ifglobal


\setcounter{secnumdepth}{4}

% >>> M5: If this module needs any updates in
% './00-shared/preamble-trm-module.sty'
% (include additional package, etc.), contact your module PM
% or create the file './00-shared/changelog.tex' make the
% necessary updates yourself and document them in this file <<<


\begin{document}


% Sets language into which pre-defined bits of text are translated
\selectlanguage{english}

% Do not delete this block
\ifglobal
% Add the following front matter if
% this module is outside of global project
\else
    \InputIfFileExists{readme}{}{}
    {\let\clearpage\relax \listoftodos}
    \tableofcontents
    %\listoftables
    %\listoffigures
\fi


%%%%%%%%%%%%%%%%%%%%%%%%%
%%% TRM Module Begins %%%
%%%%%%%%%%%%%%%%%%%%%%%%%

\hypertarget{xtwdt}{}
\chapter{XTAL32K Watchdog Timers (XTWDT)}
\label{mod:xtwdt}


\section{Overview}

The XTAL32K watchdog timer on \chipname{} is used to monitor the status of external crystal XTAL32K\_CLK. This watchdog timer can detect the oscillation failure of XTAL32K\_CLK, change the clock source of RTC, etc. When XTAL32K\_CLK works as the clock source of RTC\_SLOW\_CLK (for clock description, see Chapter \ref{mod:resclk} \textit{\nameref{mod:resclk}}) and stops oscillating, the XTAL32K watchdog timer first switches to BACKUP32K\_CLK derived from RC\_SLOW\_CLK and generates an interrupt (if the chip is in Light-sleep or Deep-sleep mode, the CPU will be woken up), and then switches back to XTAL32K\_CLK after it is restarted by software.

\begin{figure}[H]
    \centering
    \fbox{\includegraphics[width=0.8\textwidth]{\modulefiles/figures/xtal32k_watchdog}}
    \caption{XTAL32K Watchdog Timer}
    \label{fig:xtwdt}
\end{figure}

\section{Features}

\subsection{Interrupt and Wake-Up}

When the XTAL32K watchdog timer detects the oscillation failure of XTAL32K\_CLK, an oscillation failure interrupt RTC\_XTAL32K\_DEAD\_INT (for interrupt description, please refer to Chapter \ref{mod:lowpowm} \textit{\nameref{mod:lowpowm}}) is generated. At this point, the CPU will be woken up if in Light-sleep mode or Deep-sleep mode.

\subsection{BACKUP32K\_CLK}

Once the XTAL32K watchdog timer detects the oscillation failure of XTAL32K\_CLK, it replaces XTAL32K\_CLK with BACKUP32K\_CLK (with a frequency of 32 kHz or so) derived from RC\_SLOW\_CLK as RTC\_SLOW\_CLK, so as to ensure proper functioning of the system.

\section{Functional Description}

\subsection{Workflow}

\begin{enumerate}
    \item The XTAL32K watchdog timer starts counting when \hyperref[fielddesc:RTCCNTLXTAL32KWDTEN]{RTC\_CNTL\_XTAL32K\_WDT\_EN} is enabled. The counter based on RC\_SLOW\_CLK keeps counting until it detects the positive edge of XTAL\_32K and is then cleared. When the counter reaches \hyperref[fielddesc:RTCCNTLXTAL32KWDTTIMEOUT]{RTC\_CNTL\_XTAL32K\_WDT\_TIMEOUT}, it generates an interrupt or a wake-up signal and is then reset.
    \item If \hyperref[fielddesc:RTCCNTLXTAL32KAUTOBACKUP]{RTC\_CNTL\_XTAL32K\_AUTO\_BACKUP} is set and step 1 is finished, the XTAL32K watchdog timer will automatically enable BACKUP32K\_CLK as the alternative clock source of RTC\_SLOW\_CLK, to ensure the system's proper functioning and the accuracy of timers running on RTC\_SLOW\_CLK (e.g., RTC\_TIMER). For information about clock frequency configuration, please refer to Section \ref{sec:xtwdt-divisor-config}.
    \item To restore the XTAL32K watchdog timer, software restarts XTAL32K\_CLK by turning its XPD (meaning no power-down) signal off and on again via RTC\_CNTL\_XPD\_XTAL\_32K bit. Then, the XTAL32K watchdog timer switches back to XTAL32K\_CLK as the clock source of RTC\_SLOW\_CLK by clearing \hyperref[fielddesc:RTCCNTLXTAL32KWDTEN]{RTC\_CNTL\_XTAL32K\_WDT\_EN} (BACKUP32K\_CLK\_EN is also automatically cleared). If the chip is in Light-sleep or Deep-sleep mode mode, the XTAL32K watchdog timer will wake up the CPU to finish the above steps.
\end{enumerate}


\subsection{BACKUP32K\_CLK Working Principle}\label{sec:xtwdt-divisor-config}
Chips have different RC\_SLOW\_CLK frequencies due to production process variations. To ensure the accuracy of RTC\_TIMER and other timers running on RTC\_SLOW\_CLK when BACKUP32K\_CLK is at work, the divisor of BACKUP32K\_CLK should be configured according to the actual frequency of RC\_SLOW\_CLK (see details in Chapter \ref{mod:lowpowm} \textit{\nameref{mod:lowpowm}}) via RTC\_CNTL\_XTAL32K\_CLK\_FACTOR\_REG register. Each byte  in this register corresponds to a divisor component (\begin{math}x_0\end{math} \textasciitilde \begin{math}x_7\end{math}). BACKUP32K\_CLK is divided by a fraction
where the denominator is always 4, as calculated below.

\begin{gather*}
f\_back\_clk / 4 = f\_rc\_slow\_clk/S \\
S = x_0 + x_1 + ... + x_7
\end{gather*}

f\_back\_clk is the desired frequency of BACKUP32K\_CLK; f\_rc\_slow\_clk is the actual frequency of RC\_SLOW\_CLK; \begin{math}x_0\end{math} \textasciitilde \begin{math}x_7\end{math} correspond to the pulse width in high and low state of four BACKUP32K\_CLK clock signals (unit: RC\_SLOW\_CLK clock cycle).

\subsection{Configuring the Divisor Component of BACKUP32K\_CLK}

Based on principles described in Section \ref{sec:xtwdt-divisor-config}, you can configure the divisor component as follows:
\begin{itemize}
    \item Calculate the sum of divisor components S according to the frequency of RC\_SLOW\_CLK and the desired frequency of BACKUP32K\_CLK;
    \item Calculate the integer part of divisor \begin{math}N = f\_rc\_slow\_clk / f\_back\_clk\end{math};
    \item Calculate the integer part of divisor component \begin{math}M = N/2\end{math}. The integer part of divisor N are separated into two parts because a divisor component corresponds to a pulse width in high or low state;
    \item Calculate the number of divisor components that equal M (\begin{math}x_n\end{math} = M) and the number of divisor components that equal M + 1 (\begin{math}x_n\end{math} = M + 1)  according to the value of M and S. (M + 1) is the fractional part of divisor component.


\end{itemize}

For example, if the frequency of RC\_SLOW\_CLK is 163 kHz, then \begin{math}f\_rc\_slow\_clk = 163000\end{math}, \begin{math}f\_back\_clk = 32768\end{math}, \begin{math}S = 20\end{math}, \begin{math}M = 2 \end{math}, and \begin{math} \left\{x_0, x_1, x_2, x_3, x_4, x_5, x_6, x_7\right\} = \left\{2, 3, 2, 3, 2, 3, 2, 3\right\}\end{math}. As a result, the frequency of BACKUP32K\_CLK is 32.6 kHz.



\end{document}
